APPENDIX A — Mathematical Structure of the Bliss Manifold and HRD Control Law
A.1 The Balance Coefficient as a Dynamical Variable
The VIM engine is governed by the balance coefficient
\beta (t)=\frac{rAE_f(t)\, rAE_c(t)}{rAE_i(t)\, rAE_t(t)}.
All four rAE variables are time‑dependent, and each is influenced by the Harmonic Stabilizer and the HRD envelope. The Bliss State corresponds to the equilibrium condition
\beta (t)=1.
The system’s dynamics can be expressed as a set of coupled differential equations in the state vector

The stabilizer acts to regulate \mathbf{x}(t) such that \beta (t)\rightarrow 1.

A.2 The Bliss Manifold
The Bliss Manifold is defined as the set of all points in the (rAE_f,rAE_i) plane for which
\beta =1.
Solving for the manifold yields

This is a one‑dimensional curve embedded in a two‑dimensional space, but because rAE_c and rAE_t are themselves dynamical variables, the Bliss Manifold is a moving nullcline in the full four‑dimensional state space.
The manifold’s instantaneous slope is
\frac{\partial rAE_f}{\partial rAE_i}=\frac{rAE_t}{rAE_c}.
This ratio is the topological‑coherence quotient, a key invariant of the VIM system.

A.3 Harmonic Resonating Dissonance as a Driving Term
Dissonance D(t) is modeled as a structured perturbation of the vacuum manifold. Its general form is
D(t)=D_0+A(t)\sin (\omega t+\phi ),
where A(t) is a slowly varying envelope.
The effect of dissonance is to modulate the flux channel:
rAE_f(t)=rAE_f^{(0)}+\gamma _fD(t),
and to modulate impedance oppositely:
rAE_i(t)=rAE_i^{(0)}-\gamma _iD(t).
The constants \gamma _f and \gamma _i encode the system’s sensitivity to dissonance.

A.4 The HRD Control Law
The Harmonic Stabilizer applies a proportional harmonic correction to the flux and impedance channels. The discrete‑time form is


In continuous time, the stabilizer dynamics are


The signs reflect the dual nature of flux and impedance:
- When \beta <1 (leak), flux increases and impedance decreases.
- When \beta >1 (hunger), flux decreases and impedance increases.
This duality is the mathematical expression of the Theory of Balance.

A.5 Full VIM Dynamical System
Combining the HRD modulation and the stabilizer yields the full system:
\frac{d\, rAE_f}{dt}=k_f(1-\beta )+\gamma _fD(t),
\frac{d\, rAE_i}{dt}=-k_i(1-\beta )-\gamma _iD(t),
\frac{d\, rAE_c}{dt}=F_c(rAE_f,rAE_i,D(t)),
\frac{d\, rAE_t}{dt}=F_t(rAE_f,rAE_i,D(t)).
The functions F_c and F_t encode the coherence and topology responses to flux and dissonance. Their exact forms depend on the physical implementation of the rÆ‑Cell and are left intentionally open for future experimental calibration.

A.6 Stability of the Bliss Manifold
Linearizing the system around \beta =1 yields
\delta \beta =\beta -1,
and the linearized dynamics become
\frac{d\, \delta \beta }{dt}=-\lambda \delta \beta +\eta D(t),
where
\lambda =k_f\frac{\partial \beta }{\partial rAE_f}+k_i\frac{\partial \beta }{\partial rAE_i}
is the harmonic convergence rate, and
\eta =\gamma _f\frac{\partial \beta }{\partial rAE_f}-\gamma _i\frac{\partial \beta }{\partial rAE_i}
is the dissonance coupling coefficient.
The Bliss Manifold is stable when
\lambda >0.
This condition defines the allowable range of stabilizer gains (k_f,k_i).

A.7 Global Attractor Structure
Simulations confirm that the Bliss Manifold is a global attractor for a wide range of initial conditions. The attractor basin is defined by the set
\mathcal{A}=\left\{ \mathbf{x}(0)\mid \lim _{t\rightarrow \infty }\beta (t)=1\right\} .
The existence of this attractor is the mathematical justification for the rÆ‑Cell’s stability under high‑flux conditions.

A.8 Summary of Mathematical Results
- The Bliss Manifold is the nullcline \beta =1.
- HRD acts as a structured driving term that modulates flux and impedance.
- The Harmonic Stabilizer applies proportional harmonic corrections.
- The full VIM system is a four‑dimensional nonlinear dynamical system.
- Linearization shows that the Bliss Manifold is stable when \lambda >0.
- Numerical simulations confirm the existence of a global attractor.

A natural next step is to draft Appendix B — The Geometry of the rAE Alphabet, which formalizes the 24‑point Tetra‑Hexa Routing Array and its role in the VIM dynamics.
